% =============================================================================
% Project : THE_PHILOSOPHY_OF_VIABILITY
% Book    : Book I — Ontology of Viability
% Book ID : VS.B01
%
% File    : VS.B01_ONTOLOGY_OF_VIABILITY/2_MAINMATTER/PART_I_IDENTITY_AND_MANIFESTATION/CHAPTER_07_FORMAL_AND_LOGICAL_RECONSTRUCTION/sections/section_13_formal_reconstruction_of_unity_and_plurality.tex
%
% Section / Chapter:
%   Chapter 07, Section 13
%   Formal Reconstruction of Unity and Plurality
%
% Purpose :
%   To reconstruct unity and plurality through distinct formal
%   relations and structures, preserving the differences between
%   unity of an existent, unity of source, compositional unity,
%   aggregate unity, and unity of identity, as well as the
%   differences among plurality of components, relations,
%   manifestations, and identities.
%
% =============================================================================

\section{Formal Reconstruction of Unity and Plurality}
\label{sec:formal-reconstruction-of-unity-and-plurality}

\subsection{The Formal Problem of Unity and Plurality}

Unity and plurality are among the most consequential
structural notions in ontology. They concern how an
existent may constitute a unity, how multiplicity may
belong to or characterize an existent, and under what
conditions several distinguishable items may be
considered together without being reduced to an
undifferentiated aggregate.

Their formal reconstruction requires the separation
of questions that are frequently treated as if they
were interchangeable.

At least five questions must be distinguished:

\begin{enumerate}
    \item In what respect is an existent one?
    \item What kind of unity belongs to a whole or
    an organized constitution?
    \item What relation, if any, connects an existent
    with a source of unity?
    \item What forms of plurality can occur within,
    around, or in relation to an existent?
    \item Under what conditions may plurality be
    attributed to identity itself?
\end{enumerate}

These questions do not necessarily have the same
answers, and their formal representations must
not be collapsed into one predicate or relation.

The central thesis of this section is that unity
and plurality are formally distinguishable
dimensions of ontological articulation whose
specific forms must be identified through their
respective domains, respects, and relations.

A formal reconstruction must therefore represent
the relevant forms of unity and plurality without
assuming in advance that every existent possesses
exactly one identity, that every multiplicity
produces multiple identities, or that every unity
is reducible to singularity or composition.

The purpose of formalization is to make these
relations examinable. It is not to settle questions
that require independent ontological argument.

\subsection{The Domain of Formal Examination}

Let:

\[
    D_E
\]

denote a specified domain of existents.

The domain is introduced as a formal domain of
examination. Its precise extension depends on the
scope of the interpretation under consideration.

Let:

\[
    x,y,z \in D_E
\]

represent members of that domain.

The notation does not assume that every member
of \(D_E\) has the same internal constitution,
the same mode of manifestation, or the same
relation to other members.

Nor does membership in \(D_E\) establish that
the domain is an exhaustive inventory of being.

The formal reconstruction must distinguish the
domain of existents from domains of components,
relations, manifestations, and identities whenever
the chosen language requires them to be represented
separately.

A many-sorted formal language may be useful for
this purpose.

For example, the following sorts may be introduced
as provisional formal categories:

\[
    \mathsf{E}:
    \text{existents},
\]

\[
    \mathsf{C}:
    \text{components},
\]

\[
    \mathsf{R}:
    \text{relations},
\]

\[
    \mathsf{M}:
    \text{manifestations},
\]

\[
    \mathsf{I}:
    \text{identity instances or identity references}.
\]

These sorts are formal devices for controlling
the types of variables and relations. They do
not by themselves establish that the corresponding
categories are ontologically independent substances
or mutually exclusive classes.

In particular, the formal category
\(\mathsf{I}\) must not be interpreted as a
collection of properties added to existents.
Its use is intended to make claims about identity
formally expressible while preserving the
foundational commitment that being and identity
are ontologically inseparable and analytically
distinguishable.

\subsection{Unity as a Respect-Sensitive Relation}

The expression ``unity'' is incomplete when it
does not specify the respect in which something
is one.

Let:

\[
    \mathcal{U}
\]

denote a set of specified respects of unity.

For each respect \(r \in \mathcal{U}\), a predicate
may be introduced:

\[
    U_r(x).
\]

The formula:

\[
    U_r(x)
\]

states that \(x\) possesses or exhibits unity
in respect \(r\), according to the interpretation
assigned to that predicate.

The index \(r\) is essential. It prevents the
formal reconstruction from treating every
meaning of unity as a single undifferentiated
condition.

For example, the following may be distinguished:

\[
    U_{\mathrm{const}}(x),
\]

\[
    U_{\mathrm{comp}}(x),
\]

\[
    U_{\mathrm{org}}(x),
\]

\[
    U_{\mathrm{id}}(x),
\]

where the subscripts indicate, respectively,
a specified unity of constitution, composition,
organization, and identity.

These expressions are provisional formal
designations. Their precise definitions must
be established according to the relevant
ontological analysis.

The use of several unity predicates does not
presuppose that they are mutually exclusive.
An existent may satisfy more than one unity
condition.

However, the satisfaction of one condition
does not automatically entail the satisfaction
of another.

For example:

\[
    U_{\mathrm{comp}}(x)
    \not\Rightarrow
    U_{\mathrm{id}}(x)
\]

without an additional argument establishing
the relevant connection.

Likewise:

\[
    U_{\mathrm{id}}(x)
    \not\Rightarrow
    U_{\mathrm{comp}}(x)
\]

unless the meaning of identity and the relevant
compositional conditions justify that inference.

The formal advantage of a respect-sensitive
representation is that it allows the ontology
to recognize different forms of unity without
prematurely reducing them to a common structure.

\subsection{Unity of an Existent and Unity of Identity}

The unity of an existent and the unity of its
identity must be examined as distinct questions.

Let:

\[
    U_E(x)
\]

represent a specified unity attributed to an
existent \(x\).

Let:

\[
    U_I(i)
\]

represent a specified unity attributed to an
identity reference \(i\).

The distinction between these expressions
does not establish that the existent and its
identity are ontologically separable.

It distinguishes the formal questions being
examined.

The first concerns the unity of the existent
under a specified criterion. The second concerns
the unity of identity under an independently
specified account of identity.

The relationship between them requires
argument.

For example, if an existent has multiple
distinguishable components, the unity of the
existent may concern how those components
belong to or constitute it.

That fact alone does not determine whether
its identity is one or plural.

Similarly, a claim concerning the unity of
identity cannot be established merely by
showing that the existent is treated as one
in ordinary counting or reference.

The formal reconstruction therefore preserves
the distinction:

\[
    U_E(x)
    \neq
    U_I(i)
\]

as a distinction between the contents of
two claims.

It does not imply that the two forms of unity
must be unrelated. Their possible relation
remains a question for ontological analysis.

\subsection{Unity of Source}

A further form of unity concerns source.

Let:

\[
    S(u,x)
\]

represent a specified source relation between
\(u\) and \(x\).

A unity-of-source claim may be expressed
schematically as:

\[
    U_{\mathrm{source}}(x,u).
\]

This formula states that the unity of \(x\),
in the specified respect, is attributed to
or grounded in \(u\).

The precise interpretation of this relation
must determine whether it concerns origin,
constitution, dependence, derivation, or
another specified relation.

These possibilities must not be conflated.

A relation of origin does not automatically
establish a relation of present constitution.
A relation of constitution does not automatically
establish causal dependence.
A relation of dependence does not automatically
establish numerical identity between source
and dependent existent.

The formal reconstruction should therefore
use a sufficiently specific relation rather
than allowing the word ``source'' to carry
several unexamined meanings.

A claim of unity of source also differs from
a claim of unity of aggregate.

A plurality of items may be related to one
source without becoming one aggregate.
Conversely, an aggregate may be constituted
by several items without sharing a single
source in the relevant sense.

Accordingly:

\[
    U_{\mathrm{source}}(x,u)
\]

and:

\[
    U_{\mathrm{aggregate}}(x)
\]

must remain independently examinable.

Their possible connection depends on the
nature of the source relation and the
constitution of the relevant whole.

\subsection{Composition and the Unity of a Whole}

Composition concerns the constitution of
a whole through components or constituents.

Let:

\[
    C(c,x)
\]

represent a specified component relation
between \(c\) and \(x\).

A composition claim may be expressed as:

\[
    \exists c \in D_C\; C(c,x).
\]

This formula states that at least one
component is related to \(x\) under the
specified interpretation.

A plurality of components may be represented
as:

\[
    \exists c_1,c_2 \in D_C
    \bigl(
        C(c_1,x)
        \land
        C(c_2,x)
        \land
        c_1 \neq c_2
    \bigr).
\]

This establishes the presence of at least
two formally distinguishable components
related to \(x\).

It does not, by itself, establish the
conditions under which those components
constitute one whole.

A stronger compositional account would
require additional relations or conditions,
such as:

\begin{itemize}
    \item the criterion by which components
    belong to the whole;
    \item the relation among the components;
    \item the conditions under which the
    whole is constituted;
    \item the relevant limits of the
    component relation;
    \item the conditions under which
    component replacement or alteration
    affects the whole.
\end{itemize}

The formal reconstruction must not assume
that composition is always additive or
that a whole is exhausted by a set of
components.

The relation between a whole and its
components may require a more complex
formal structure than simple set membership.

In particular, the expression:

\[
    x = \{c_1,c_2,\ldots,c_n\}
\]

should not be adopted as a general
definition of a composed existent.

It would identify the existent with a
set of components and would thereby
introduce a substantial ontological
commitment requiring independent
justification.

Composition must therefore be represented
through its own relations and criteria.

\subsection{Aggregate and Compositional Unity}

An aggregate is a plurality considered
together under a specified criterion.

Let:

\[
    A(a,\Gamma)
\]

represent an aggregate \(a\) associated
with a collection \(\Gamma\).

The collection may be formally represented
as a set, sequence, or other suitable
structure, depending on the intended
interpretation.

The aggregate relation is not identical
to the relation of composition.

A collection of items may be treated as
an aggregate without the ontology having
established that the collection constitutes
a unified existent in a stronger sense.

Conversely, a composed whole may possess
a unity that cannot be adequately represented
by treating it as a mere collection.

The distinction may be expressed as:

\[
    \text{Aggregate Unity}
    \neq
    \text{Compositional Unity}.
\]

This distinction is not a denial that
some aggregates may also be composed
wholes.

It means that the two descriptions
require different criteria.

For an aggregate, the relevant question
may be what makes the items count together.

For a composed whole, the question
concerns the relation through which
the components constitute that whole.

A formal system must therefore specify
whether a claim concerns collection,
constitution, organization, or another
form of unity.

Without that specification, the expression
``one whole'' may conceal several distinct
ontological commitments.

\subsection{Singularity and Unity}

Singularity concerns a specified mode
of being one, whereas unity concerns
a form or condition of unification.

The two notions may be related, but
their equivalence cannot be assumed.

Let:

\[
    \operatorname{Sing}(x)
\]

represent a specified singularity
condition.

Let:

\[
    U_r(x)
\]

represent unity in respect \(r\).

The formal language permits the
independent examination of:

\[
    \operatorname{Sing}(x),
\]

\[
    U_r(x),
\]

and the possible relations between them.

For example, a particular interpretation
may establish:

\[
    \operatorname{Sing}(x)
    \rightarrow U_r(x).
\]

Another interpretation may require
additional conditions before such an
implication is justified.

The ontology must not assume a universal
biconditional:

\[
    \operatorname{Sing}(x)
    \leftrightarrow U_r(x).
\]

The validity of any such formula depends
on the meanings assigned to singularity
and the relevant respect of unity.

This distinction is especially important
when considering plurality. A structure
may exhibit unity while containing
multiple distinguishable constituents.
The presence of those constituents does
not automatically eliminate unity.

Likewise, the fact that something is
counted as one does not settle the
nature of its unity.

The formal reconstruction must therefore
keep singularity and unity available
as distinct analytical categories.

\subsection{Plurality as a Typed Formal Relation}

Plurality must be specified according
to what is plural and the respect in
which multiplicity is asserted.

Let:

\[
    \mathcal{P}
\]

denote a set of specified plurality
types.

For each type \(p \in \mathcal{P}\),
let:

\[
    \operatorname{Pl}_p(x)
\]

represent the corresponding plurality
condition associated with \(x\).

The following distinctions are
particularly important:

\begin{itemize}
    \item plurality of components;
    \item plurality of relations;
    \item plurality of manifestations;
    \item plurality of determinations;
    \item plurality of identity instances;
    \item plurality of existents.
\end{itemize}

These forms of plurality have different
formal subjects and different conditions
of attribution.

For example, plurality of components
may be represented through the component
relation \(C\).

Plurality of relations may require a
domain of relations and a criterion
for distinguishing relational instances.

Plurality of manifestations may require
a domain of manifestation instances
and a relation connecting them to an
existent.

Plurality of identities requires an
independent account of what counts as
a distinct identity and under which
respect or level the distinction is
being asserted.

The formal system must not use one
plurality predicate as if it established
all these forms simultaneously.

\subsection{Plurality of Components}

Plurality of components can be expressed
through the existence of distinguishable
component instances:

\[
    \exists c_1,c_2 \in D_C
    \bigl(
        C(c_1,x)
        \land
        C(c_2,x)
        \land
        c_1 \neq c_2
    \bigr).
\]

The formula establishes component plurality
under the selected interpretation.

It does not determine:

\begin{enumerate}
    \item whether the components are
    independently existent;
    \item whether they are mutually
    dependent;
    \item whether they constitute one
    whole;
    \item whether they belong to different
    organizational levels;
    \item whether their plurality affects
    the identity of the existent;
    \item whether the components themselves
    possess identities.
\end{enumerate}

Each of these questions requires
additional formal and ontological
conditions.

The relation \(C(c,x)\) must also be
defined carefully. A component may be
a material constituent, a structural
constituent, or another type of
constituent. The formal system should
not assume that all component relations
have the same ontological meaning.

A more refined language may therefore
use typed component relations:

\[
    C_{\mathrm{mat}}(c,x),
\]

\[
    C_{\mathrm{str}}(c,x),
\]

\[
    C_{\mathrm{other}}(c,x).
\]

These expressions remain provisional
until the relevant types have been
defined in the ontology.

Their purpose is to prevent the
formalization from treating every
distinguishable constituent as
ontologically equivalent.

\subsection{Plurality of Relations}

An existent may stand in multiple
relations to other existents or to
other elements of its ontological
context.

Let:

\[
    R(r,x,y)
\]

represent a relation instance \(r\)
connecting \(x\) and \(y\).

Plurality of relations may be expressed
as:

\[
    \exists r_1,r_2 \in D_R
    \bigl(
        R(r_1,x,y)
        \land
        R(r_2,x,y)
        \land
        r_1 \neq r_2
    \bigr).
\]

This formula states that two
distinguishable relation instances
connect the relevant relata under
the chosen interpretation.

The existence of several relations
does not automatically establish
several existents or several identities.

Nor does it establish that the
relations are independent of the
relata.

The formal system must distinguish
the plurality of relation instances
from the plurality of relata and
from the plurality of identity
instances.

A relationally complex existent
may therefore be examined without
assuming that every relational
difference constitutes a difference
of identity.

\subsection{Plurality of Manifestations}

Plurality of manifestations concerns
the occurrence of distinguishable
manifestation instances associated
with an existent.

Let:

\[
    M(m,x)
\]

represent a specified manifestation
relation between manifestation instance
\(m\) and existent \(x\).

Plurality of manifestations may be
expressed as:

\[
    \exists m_1,m_2 \in D_M
    \bigl(
        M(m_1,x)
        \land
        M(m_2,x)
        \land
        m_1 \neq m_2
    \bigr).
\]

This formula represents distinguishable
manifestation instances associated
with the same existent.

The exact meaning of distinction
between manifestations must be
specified.

It may concern different modes,
conditions, contexts, configurations,
or other relevant respects.

The formal expression does not
determine whether the manifestations
are simultaneous, successive,
dependent, or mutually exclusive.

Those questions require additional
relations and temporal or structural
conditions where appropriate.

Most importantly, plurality of
manifestations does not by itself
establish plurality of identities.

The inference:

\[
    \operatorname{Pl}_{M}(x)
    \rightarrow
    \operatorname{Pl}_{I}(x)
\]

is not valid without an independently
justified principle connecting the
relevant manifestation distinctions
to identity distinctions.

The formal reconstruction must
therefore preserve the possibility
that one existent may exhibit
multiple manifestations while the
question of the number of its
identities remains separately open.

\subsection{Plurality of Identities}

Plurality of identities is a distinct
ontological problem.

It concerns whether an existent may
possess or be associated with more
than one identity in a specified
respect, level, or ontological
arrangement.

Let:

\[
    I(i,x)
\]

represent a specified relation
between identity reference \(i\)
and existent \(x\).

A formal expression of identity
plurality may be written:

\[
    \exists i_1,i_2 \in D_I
    \bigl(
        I(i_1,x)
        \land
        I(i_2,x)
        \land
        i_1 \neq i_2
    \bigr).
\]

This expression is only a schema.
Its interpretation depends on an
adequate account of identity references,
the relation \(I\), and the criterion
by which two identity references
count as distinct.

The formula must not be interpreted
as a claim that identities are
properties attached to an existent.

It is a formal representation of
a question concerning the possible
multiplicity of identity.

The distinctness condition:

\[
    i_1 \neq i_2
\]

is a formal condition within the
chosen language. It does not
automatically establish a substantive
ontological difference unless the
interpretation of identity references
is justified.

The formal reconstruction must
therefore distinguish:

\begin{itemize}
    \item the existence of multiple
    identity references in a model;
    \item the formal distinctness of
    those references;
    \item the ontological distinctness
    of the identities represented;
    \item the relation of those identities
    to one existent;
    \item the respect or level under
    which their plurality is asserted.
\end{itemize}

The possibility of identity plurality
remains an open ontological question
whose resolution cannot be obtained
merely by counting components,
relations, or manifestations.

\subsection{Respect and Level in Identity Plurality}

An adequate analysis of identity
plurality must specify the respect
and level under which identity
distinctions are proposed.

Let:

\[
    \rho
\]

denote a specified respect, and let:

\[
    \lambda
\]

denote a specified level of analysis.

A refined identity relation may
therefore be represented as:

\[
    I(i,x,\rho,\lambda).
\]

The additional parameters make it
possible to distinguish claims that
would otherwise be treated as
equivalent.

For example, two identity references
may be formally distinguishable
under different respects without
thereby establishing two identities
at the same level.

Likewise, two identity references
associated with different levels
must not automatically be treated
as instances of a hierarchical
identity structure.

The introduction of levels does
not itself establish:

\begin{itemize}
    \item a hierarchy of identities;
    \item a higher identity governing
    lower identities;
    \item a linear order among identity
    levels;
    \item a universal rule of identity
    aggregation;
    \item a fixed number of identity
    levels.
\end{itemize}

Each such structure would require
additional principles and arguments.

The purpose of the parameters
\(\rho\) and \(\lambda\) is to
make the conditions of an identity
plurality claim explicit.

They do not predetermine the result
of the ontological inquiry.

\subsection{The Non-Transfer of Plurality Between Domains}

One of the principal requirements
of the formal reconstruction is
to prevent an inference from one
plurality domain to another without
an explicit bridge principle.

Suppose:

\[
    \operatorname{Pl}_{C}(x)
\]

expresses plurality of components,

\[
    \operatorname{Pl}_{R}(x)
\]

expresses plurality of relations,

\[
    \operatorname{Pl}_{M}(x)
\]

expresses plurality of manifestations,
and:

\[
    \operatorname{Pl}_{I}(x)
\]

expresses plurality of identities.

The formal system must not assume
the unrestricted implications:

\[
    \operatorname{Pl}_{C}(x)
    \rightarrow
    \operatorname{Pl}_{I}(x),
\]

\[
    \operatorname{Pl}_{R}(x)
    \rightarrow
    \operatorname{Pl}_{I}(x),
\]

\[
    \operatorname{Pl}_{M}(x)
    \rightarrow
    \operatorname{Pl}_{I}(x).
\]

Each implication would require
an independent argument identifying
the conditions under which the
relevant plurality entails identity
plurality.

The same restriction applies
in the opposite direction.

A claim of identity plurality
does not automatically establish
component, relational, or
manifestational plurality.

The formal distinction is therefore
not merely terminological. It
prevents the transfer of conclusions
between different ontological
domains without justified premises.

\subsection{Unity and Plurality Within One Formal Structure}

Unity and plurality may occur
together when they concern different
respects or formal relations.

A structure may satisfy:

\[
    U_r(x)
\]

while also satisfying:

\[
    \operatorname{Pl}_{p}(x).
\]

This conjunction is formally
consistent unless the definitions
or additional axioms of the system
exclude it.

For example, a composed whole may
possess a specified unity while
containing multiple distinguishable
components.

The formal representation is:

\[
    U_{\mathrm{comp}}(x)
    \land
    \operatorname{Pl}_{C}(x).
\]

This conjunction does not generate
a contradiction merely because
unity and plurality occur together.

The apparent conflict arises only
when unity and plurality are
treated as claims about precisely
the same respect under mutually
exclusive conditions.

A formal reconstruction must
therefore identify the respect
of unity and the type of plurality
before testing their compatibility.

This allows the ontology to
investigate structures in which
unity and plurality coexist
without reducing one to the other.

\subsection{Unity, Plurality, and the Principles of Being and Identity}

The formal reconstruction must
remain compatible with the three
foundational principles concerning
being and identity.

First, Analytical Distinction
requires that the questions
concerning being and identity
remain analytically distinguishable.

Accordingly, a formal unity
predicate must not be used as
a substitute for the identity
relation itself.

Second, Ontological Inseparability
requires that being and identity
remain ontologically inseparable
within the adopted foundational
framework.

The formal representation of
unity and plurality must therefore
not introduce an interpretation
in which an existent is treated
as ontologically complete while
identity is merely an optional
external attribute.

Third, Absence of Ontological
Priority requires that the
formal reconstruction avoid
establishing an unsupported
priority of being over identity
or identity over being.

These principles constrain the
interpretation of the formal
vocabulary, but they do not
automatically determine whether
unity or plurality should be
attributed to identity in a
particular case.

The question of identity plurality
must therefore be examined
without weakening the foundational
principles or using them as
substitutes for the required
argument.

\subsection{Axioms, Definitions, and Propositions Concerning Unity and Plurality}

The formal reconstruction requires
a clear distinction between
stipulative definitions, candidate
axioms, and propositions whose
validity depends on further
premises.

The following definitions are
proposed as formal conventions
for the present analysis.

\paragraph{Definition 13.1: Respect-sensitive unity.}

For a specified respect \(r\),
\(U_r(x)\) expresses that \(x\)
satisfies the criterion of unity
assigned to \(r\).

\paragraph{Definition 13.2: Typed plurality.}

For a specified plurality type
\(p\), \(\operatorname{Pl}_p(x)\)
expresses that \(x\) satisfies
the criterion of multiplicity
assigned to \(p\).

\paragraph{Definition 13.3: Component plurality.}

\(\operatorname{Pl}_{C}(x)\)
expresses that at least two
formally distinguishable
components are related to \(x\)
under the specified component
relation.

\paragraph{Definition 13.4: Manifestation plurality.}

\(\operatorname{Pl}_{M}(x)\)
expresses that at least two
formally distinguishable
manifestation instances are
related to \(x\) under the
specified manifestation relation.

\paragraph{Definition 13.5: Identity plurality.}

\(\operatorname{Pl}_{I}(x)\)
expresses that at least two
identity references are related
to \(x\) under a specified
interpretation and satisfy
the criterion of identity
distinctness adopted by the
formal system.

These definitions specify the
formal vocabulary. They do not
establish the ontological truth
of every claim expressed through
that vocabulary.

The following propositions
distinguish results that can
be maintained at the current
level of formal reconstruction.

\paragraph{Proposition 13.1: Unity requires a specified respect.}

An unqualified unity claim is
formally incomplete whenever
different respects of unity
remain relevant to its intended
meaning.

\paragraph{Proposition 13.2: Unity and plurality are not universally exclusive.}

A unity condition and a plurality
condition may be jointly satisfied
when they concern compatible
respects or types.

\paragraph{Proposition 13.3: Component plurality does not entail identity plurality.}

The existence of multiple
distinguishable components does
not, without an additional
bridge principle, entail multiple
identities.

\paragraph{Proposition 13.4: Manifestation plurality does not entail identity plurality.}

The existence of multiple
manifestation instances does
not, without an additional
bridge principle, establish
multiple identities.

\paragraph{Proposition 13.5: Aggregate unity and compositional unity require distinct criteria.}

The formal representation of
an aggregate does not by itself
establish the compositional
unity of an existent.

\paragraph{Proposition 13.6: Unity and singularity require independent analysis.}

A claim concerning unity cannot
be substituted for a claim
concerning singularity without
an argument establishing the
relevant equivalence or implication.

\paragraph{Proposition 13.7: Identity plurality remains an independent ontological question.}

The formal possibility of
representing multiple identity
references does not settle
whether an existent is
ontologically characterized
by multiple identities.

These propositions have different
justificatory statuses.

Some express consequences of
the formal distinctions introduced
here. Others formulate methodological
constraints on valid inference.
The ontological truth of claims
concerning identity plurality
remains open pending further
argument.

\subsection{Countermodels and the Examination of Unwarranted Inference}

The model-based method introduced
in Section 11 can be used to
test proposed implications
between different forms of
plurality.

Consider the proposed inference:

\[
    \operatorname{Pl}_{C}(x)
    \rightarrow
    \operatorname{Pl}_{I}(x).
\]

A countermodel to this implication
would satisfy component plurality
while failing the specified
identity-plurality condition.

Such a model would demonstrate
that the implication is not
entailed by the formal definitions
and premises alone.

It would not establish that
identity plurality is impossible
in every ontological case.

It would establish only that
component plurality, as formally
specified, is insufficient to
entail identity plurality.

The same procedure can be used
for proposed implications between
manifestation plurality and
identity plurality, aggregate
unity and compositional unity,
or unity and singularity.

Every countermodel must identify
the exact formal assumptions
under examination.

The result must not be transferred
beyond those assumptions.

\subsection{The Limits of the Formal Reconstruction}

The formal reconstruction
establishes a framework for
distinguishing and examining
different forms of unity and
plurality.

It does not yet provide a
complete ontology of every
possible unity relation or
plurality type.

Several issues require further
development:

\begin{enumerate}
    \item the precise ontological
    conditions of unity of
    constitution;

    \item the relation between
    source unity and the unity
    of a constituted existent;

    \item the conditions under
    which composition constitutes
    a genuine whole;

    \item the relation between
    aggregate unity and the
    ontological unity of an
    existent;

    \item the possible forms
    of plurality of identity;

    \item the role of respect
    and level in distinguishing
    identity instances;

    \item the possibility of
    identity plurality without
    a corresponding plurality
    of components;

    \item the possibility of
    multiple manifestations
    without a change in the
    identity structure;

    \item the formal conditions
    under which unity and
    plurality may coexist.
\end{enumerate}

These questions must be developed
through further ontological
argument and, where appropriate,
through more precise formal
languages.

The introduction of a formal
distinction does not itself
resolve the philosophical
problem represented by that
distinction.

\subsection{Open Problems}

The following problems remain
open for subsequent examination.

First, it remains necessary
to determine whether unity
has a common ontological
foundation across its different
respects or whether the term
designates several irreducible
forms of unification.

Second, the relation between
unity of source and unity of
identity requires further
analysis.

Third, the conditions under
which composition constitutes
a genuine unity must be
distinguished from the
conditions under which
components are merely
collected or grouped.

Fourth, the ontology must
examine whether identity
plurality is possible, and
if so, under what respects
and levels it may occur.

Fifth, it remains necessary
to determine whether the
plurality of identities,
where established, has any
necessary relation to the
plurality of existents.

Sixth, the formal treatment
of unity must be examined
in relation to the concepts
of expression and manifestation
developed elsewhere in Part I.

Seventh, the possibility
of forms of unity that are
not adequately represented
by ordinary set-theoretical
or mereological structures
requires further investigation.

These problems remain open
because their resolution
requires more than the
formal vocabulary introduced
in this section.

\subsection{Concluding Principle}

Unity and plurality are not
single formal conditions
opposed to one another.
They are families of
ontologically significant
relations and structures
whose meanings depend on
the respect, domain, and
type under examination.

A formal reconstruction
must distinguish unity of
an existent, unity of
composition, aggregate
unity, unity of source,
and unity of identity.
It must likewise distinguish
plurality of components,
relations, manifestations,
and identities.

The governing principle
may therefore be stated
as follows:

\begin{quote}
A claim of unity or plurality
is ontologically determinate
only when the subject,
respect, and type of the
claim have been specified.
No form of plurality
automatically establishes
another form of plurality,
and no form of unity
automatically determines
the unity or singularity
of identity.
\end{quote}

This formal reconstruction
preserves the possibility
of unity and plurality
within the same ontological
structure while preventing
unjustified transfers of
meaning and inference
between distinct domains.

The next section examines
the formal reconstruction
of expression and
manifestation, focusing
on the relations through
which an existent may
be expressed or manifested
without reducing either
notion to identity itself.

% =============================================================================
% END OF SECTION 13 OF CHAPTER 07
% =============================================================================